% GENERATED FILE. Edit canonical lessons, then run npm run generate:books.
% canonical-source-hash: fnv1a64:7ba2936699f50187
% canonical-lessons: KA-C277-uhisu, KA-C277-anisu, KA-C277-abhyasa-madu, KA-C277-banna-haccu, KA-C277-daunlod-madu

\chapter{To guess, To seem, To practise, To paint, To download}
\label{ch:ka-to-guess-to-seem-to-practise-to-paint-to-download}
\hlchaptermodality{277}

\begin{chapteropening}
\textbf{By the end of this chapter:} \emph{I can say to guess, to seem, to practise, to paint and to download.}

Complete the last lesson of chapter 277: I can say to guess, to seem, to practise, to paint and to download.
\end{chapteropening}

\section[\texorpdfstring{\kn{ಊಹಿಸು} (ūhisu)}{ಊಹಿಸು (ūhisu)}]{\kn{ಊಹಿಸು} (ūhisu) — to guess}

\label{lesson:KA-C277-uhisu}



\begin{quote}
Before the new one: say the Kannada for to console, then the Kannada for to be proud.
\end{quote}

\subsection*{You'll want to know: \kn{ಊಹಿಸು}}
\textbf{\kn{ಊಹಿಸು}} — \emph{ūhisu} — \textquotedblleft{}to guess\textquotedblright{}.

\begin{grammarlens}[title={What you've built}]
One more word for this chapter.
\end{grammarlens}

\subsection*{Your turn}
\begin{itemize}
\raggedright
  \item \emph{Say it:} \emph{ūhisu}
  \item \emph{Say it:} \emph{ūhisu}, once more
  \item \emph{Say it:} say \emph{ūhisu}
  \item \emph{Recall:} say the Kannada for to console, then the Kannada for to be proud, then say \emph{ūhisu} again
\end{itemize}

\begin{tcolorbox}[breakable,colback=teal!4,colframe=teal!35!black,title={Before you move on}]
What does \kn{ಊಹಿಸು} mean? (\textquotedblleft{}To guess\textquotedblright{}.) Say it once more.
\end{tcolorbox}
\section[\texorpdfstring{\kn{ಅನಿಸು} (anisu)}{ಅನಿಸು (anisu)}]{\kn{ಅನಿಸು} (anisu) — to seem, to feel (that)}

\label{lesson:KA-C277-anisu}



\begin{quote}
Before the new one: say the Kannada for to be proud, then the Kannada for to guess.
\end{quote}

\subsection*{You'll want to know: \kn{ಅನಿಸು}}
\textbf{\kn{ಅನಿಸು}} — \emph{anisu} — \textquotedblleft{}to seem, to feel (that)\textquotedblright{}.

\begin{grammarlens}[title={What you've built}]
One more word for this chapter.
\end{grammarlens}

\subsection*{Your turn}
\begin{itemize}
\raggedright
  \item \emph{Say it:} \emph{anisu}
  \item \emph{Say it:} \emph{anisu}, once more
  \item \emph{Say it:} say \emph{anisu}
  \item \emph{Recall:} say the Kannada for to be proud, then the Kannada for to guess, then say \emph{anisu} again
\end{itemize}

\begin{tcolorbox}[breakable,colback=teal!4,colframe=teal!35!black,title={Before you move on}]
What does \kn{ಅನಿಸು} mean? (\textquotedblleft{}To seem, to feel (that)\textquotedblright{}.) Say it once more.
\end{tcolorbox}
\section[\texorpdfstring{\kn{ಅಭ್ಯಾಸ} \kn{ಮಾಡು} (abhyāsa māḍu)}{ಅಭ್ಯಾಸ ಮಾಡು (abhyāsa māḍu)}]{\kn{ಅಭ್ಯಾಸ} \kn{ಮಾಡು} (abhyāsa māḍu) — to practise}

\label{lesson:KA-C277-abhyasa-madu}



\begin{quote}
Before the new one: say the Kannada for to guess, then the Kannada for to seem.
\end{quote}

\subsection*{You'll want to know: \kn{ಅಭ್ಯಾಸ} \kn{ಮಾಡು}}
\textbf{\kn{ಅಭ್ಯಾಸ} \kn{ಮಾಡು}} — \emph{abhyāsa māḍu} — \textquotedblleft{}to practise\textquotedblright{}.

\begin{grammarlens}[title={What you've built}]
One more word for this chapter.
\end{grammarlens}

\subsection*{Your turn}
\begin{itemize}
\raggedright
  \item \emph{Say it:} \emph{abhyāsa māḍu}
  \item \emph{Say it:} \emph{abhyāsa māḍu}, once more
  \item \emph{Say it:} say \emph{abhyāsa māḍu}
  \item \emph{Recall:} say the Kannada for to guess, then the Kannada for to seem, then say \emph{abhyāsa māḍu} again
\end{itemize}

\begin{tcolorbox}[breakable,colback=teal!4,colframe=teal!35!black,title={Before you move on}]
What does \kn{ಅಭ್ಯಾಸ} \kn{ಮಾಡು} mean? (\textquotedblleft{}To practise\textquotedblright{}.) Say it once more.
\end{tcolorbox}
\section[\texorpdfstring{\kn{ಬಣ್ಣ} \kn{ಹಚ್ಚು} (baṇṇa haccu)}{ಬಣ್ಣ ಹಚ್ಚು (baṇṇa haccu)}]{\kn{ಬಣ್ಣ} \kn{ಹಚ್ಚು} (baṇṇa haccu) — to paint (apply colour)}

\label{lesson:KA-C277-banna-haccu}



\begin{quote}
Before the new one: say the Kannada for to seem, then the Kannada for to practise.
\end{quote}

\subsection*{You'll want to know: \kn{ಬಣ್ಣ} \kn{ಹಚ್ಚು}}
\textbf{\kn{ಬಣ್ಣ} \kn{ಹಚ್ಚು}} — \emph{baṇṇa haccu} — \textquotedblleft{}to paint (apply colour)\textquotedblright{}.

\begin{grammarlens}[title={What you've built}]
One more word for this chapter.
\end{grammarlens}

\subsection*{Your turn}
\begin{itemize}
\raggedright
  \item \emph{Say it:} \emph{baṇṇa haccu}
  \item \emph{Say it:} \emph{baṇṇa haccu}, once more
  \item \emph{Say it:} say \emph{baṇṇa haccu}
  \item \emph{Recall:} say the Kannada for to seem, then the Kannada for to practise, then say \emph{baṇṇa haccu} again
\end{itemize}

\begin{tcolorbox}[breakable,colback=teal!4,colframe=teal!35!black,title={Before you move on}]
What does \kn{ಬಣ್ಣ} \kn{ಹಚ್ಚು} mean? (\textquotedblleft{}To paint (apply colour)\textquotedblright{}.) Say it once more.
\end{tcolorbox}
\section[\texorpdfstring{\kn{ಡೌನ್ಲೋಡ್} \kn{ಮಾಡು} (ḍaunlōḍ māḍu)}{ಡೌನ್ಲೋಡ್ ಮಾಡು (ḍaunlōḍ māḍu)}]{\kn{ಡೌನ್ಲೋಡ್} \kn{ಮಾಡು} (ḍaunlōḍ māḍu) — to download}

\label{lesson:KA-C277-daunlod-madu}



\begin{quote}
Before the new one: say the Kannada for to practise, then the Kannada for to paint.
\end{quote}

\subsection*{You'll want to know: \kn{ಡೌನ್ಲೋಡ್} \kn{ಮಾಡು}}
\textbf{\kn{ಡೌನ್ಲೋಡ್} \kn{ಮಾಡು}} — \emph{ḍaunlōḍ māḍu} — \textquotedblleft{}to download\textquotedblright{}.

\begin{grammarlens}[title={What you've built}]
One more word for this chapter.
\end{grammarlens}

\subsection*{Your turn}
\begin{itemize}
\raggedright
  \item \emph{Say it:} \emph{ḍaunlōḍ māḍu}
  \item \emph{Say it:} \emph{ḍaunlōḍ māḍu}, once more
  \item \emph{Say it:} say \emph{ḍaunlōḍ māḍu}
  \item \emph{Recall:} say the Kannada for to practise, then the Kannada for to paint, then say \emph{ḍaunlōḍ māḍu} again
\end{itemize}

\begin{tcolorbox}[breakable,colback=teal!4,colframe=teal!35!black,title={Before you move on}]
What does \kn{ಡೌನ್ಲೋಡ್} \kn{ಮಾಡು} mean? (\textquotedblleft{}To download\textquotedblright{}.) Say it once more.
\end{tcolorbox}
